\section{Discussion and limitations}\label{sec:discussion}
\paragraph{What the evidence says.}
Under our model the answer to ``what investment puts a target first, and how does it vary with the rule'' is conditional on three things. \emph{The cost measure}: in total points the mean is among the dearest rules, and the minimum (on both leaderboards) and the median (on HELM) are several times cheaper, but not in the largest single gain, and in tasks touched the order is different and not stable (Sections~\ref{sec:headline}, \ref{sec:order}). \emph{The cap}: whether first place is reachable at all under bounded per-task gains depends on the rule, and on OLLB for most targets the answer is no at 10 points under every rule (Section~\ref{sec:caps}). \emph{The target's position and the noise}: targets near the top need few points, but the nominal minimum succeeds in a median of only \num{r2.ollb.obs1.p0min} to \num{r2.ollb.obs1.p0max} (OLLB) and \num{r2.helm.obs1.p0min} to \num{r2.helm.obs1.p0max} (HELM) of noise draws (Section~\ref{sec:noise}). Of the six pre-stated claims, C1 holds numerically, C2 is refuted as a general order, C3 is supported, C4 holds at the ceiling and is qualified by feasibility, C5 holds in its probabilistic form, and C6 compares curves with and without correlation.

\paragraph{Where the conclusions fail or do not apply.}
(i) \emph{Cost is in leaderboard points}, not in data, compute or money; the caps stand in for what is plausible, and the empirical caps (\textsf{emp95}, \textsf{emp99}) are general fine-tuning gains on six tasks, not the effect of task-specific data. (ii) \emph{Rivals are fixed} and tasks do not interact: no spillover of a gain to other tasks (which would lower the cost) and no response from other developers (which would raise it; the equilibrium question is studied by \citet{chen2026leaderboard}). (iii) \emph{Two leaderboards}, both narrow: six heterogeneously normalised tasks on OLLB and one benchmark split into 57 subjects on HELM; the $n$-sweep varies $n$ only within MMLU. (iv) \emph{Rank-rule coverage}: the Condorcet requirement stands for Condorcet-consistent rules, not for Kemeny (not run); on OLLB rank rules run on a rank-stratified sample whose quantiles are not those of all targets. (v) \emph{Normalisation}: OLLB scores are the leaderboard's normalised points; a raw-score slice keeps the order of rules but not the magnitudes. (vi) \emph{Ties and discreteness}: $\varepsilon$ is small and its effect bounded (Section~\ref{sec:robust}). (vii) \emph{Dependence}: intervals on target medians describe the set of targets studied, not independent samples of models and not evaluation uncertainty; bootstrapping OLLB over \num{cb.nclusters} base-model families widens them by a median factor of \num{cb.width.med} and changes no conclusion (Section~\ref{sec:caps}), while HELM has no family labels. (viii) \emph{Noise model}: Gaussian with item-sampling errors only; the overdispersion factors are sensitivity multipliers, not estimates, and the regression-to-the-mean reading rests on stated assumptions.

\paragraph{Unfinished or deviating work.}
The OLLB sample is a deterministic grid, so no seed applied to it; the first noise and bootstrap runs used the script default seed 0, not the study seed, and re-runs with the study seed and redrawn samples agree within the stated intervals. HELM rank-rule noise runs use 100 instead of 200 draws (exact optimisation cost), and the total compute exceeded the pre-run estimate of 4 CPU hours (Appendices~\ref{app:repro}, \ref{app:newanalyses}).

\paragraph{Implication.}
On these two leaderboards, and under our assumptions, a designer who worries about targeted investment should not infer robustness from the rule's name: whether the median or the minimum is harder to game than the mean depends on whether the developer pays in points or in tasks, and on how far a single score can plausibly move. Reporting a rule with an assumed largest gain per task, and the share of models that could reach first place within it, is a cheap check; the assumed gain itself is not something this study supplies.
